% !PRINTED 671
\weberpage{671}
Let $\Omega$ then be a field of degree $n$, and

\weq{1}{\Omega_1,\ \Omega_2,\ \ldots,\ \Omega_n}

the conjugate fields. The irreducible equation of degree $n$ whose roots determine these conjugate fields for us will in general have both real and imaginary roots; the latter, however, always occur in pairs, so that to every imaginary field $\Omega^{(i)}$ there belongs another whose numbers are conjugate imaginary to those of $\Omega^{(i)}$, and which therefore arises from $\Omega^{(i)}$ by the substitution $(i, -i)$.

Such imaginary pairs we combine into a unit, and we denote \gesperrt{the number of real fields and of pairs of imaginary fields of the sequence (1)} by $\nu$. If all the conjugate fields are real, then $n = \nu$; if they are all imaginary, then $n = 2\nu$.

Apart from the case of the rational field, $\nu$ is $= 1$ only in the one further case of an imaginary quadratic field. Otherwise $\nu$ is always greater than 1.

The number of imaginary pairs is $n - \nu$, and the number of real fields $2\nu - n$.

In the determinant by which the square root of the fundamental number, $\sqrt{\Delta}$, was defined in \S.\,145, the substitution $(i, -i)$ interchanges each two conjugate imaginary rows, and $\sqrt{\Delta}$ therefore changes its sign whenever this number is odd, but not otherwise. In the first case $\sqrt{\Delta}$ is imaginary, in the second real. From this it follows that the \gesperrt{fundamental number $\Delta$ of the field $\Omega$ is positive or negative according as $n - \nu$ is even or odd.}

If of two conjugate imaginary fields we retain only the one, there results the sequence of conjugate fields

\[ \Omega_1,\ \Omega_2,\ \ldots,\ \Omega_\nu, \]

to which we assign a system of symbols

\[ \delta_1,\ \delta_2,\ \ldots,\ \delta_\nu \]

with the provision that $\delta_s = 1$ is to hold when $\Omega_s$ is real, and $= 2$ when $\Omega_s$ is imaginary, so that $\Sigma \delta_s = n$.

Let $\eta$ denote a number of the field $\Omega$ different from zero, and

\[ \eta_1,\ \eta_2,\ \ldots,\ \eta_n \]

the conjugate numbers. To this system of numbers we assign another system of numbers

\[ \lambda_1,\ \lambda_2,\ \ldots,\ \lambda_n, \]
%CONT>>
